\chapter{TỔNG QUAN BỘ DỮ LIỆU}

\section{Nguồn gốc và cách sử dụng dữ liệu}

Các dữ liệu trong đề tài đều do nhóm mô phỏng bằng Python, không lấy từ khách hàng hay hệ thống ngân hàng thật. Nhóm dùng dữ liệu để minh họa DBSCAN trong không gian 2D, chạy các thao tác thêm/xóa và benchmark trên một luồng nhiều chiều.

Dữ liệu chính gồm tập tọa độ ATM 2D và luồng giao dịch thẻ 5D. Ngoài ra, nhóm chuẩn bị tám ví dụ nhỏ để minh họa các trường hợp cập nhật và một số ca nhỏ để kiểm thử tình huống biên.

\section{Mô tả các bộ dữ liệu}

\begin{table}[H]
\centering
\caption{Các bộ dữ liệu được dùng trong đề tài}
\label{tab:dataset_overview}
\normalsize
\begin{tabularx}{\textwidth}{@{}YP{4cm}P{2.5cm}@{}}
\toprule
\textbf{Bộ dữ liệu} & \textbf{Quy mô} & \textbf{Số chiều} \\
\midrule
ATM 2D & 1.200 điểm & 2 \\
Luồng thẻ & 500 điểm ban đầu, 300 cập nhật & 5 thuộc tính gốc \\
Ví dụ cập nhật & 8 kịch bản & 2 \\
Ca kiểm thử bổ sung & Một số ca nhỏ & 1D/2D \\
\bottomrule
\end{tabularx}
\end{table}

\textbf{ATM 2D.} Tập này có 1.200 điểm, mỗi điểm gồm hai tọa độ \texttt{x\_km} và \texttt{y\_km}. Các tọa độ mô phỏng vị trí trong mặt phẳng, đơn vị kilômét. Cột \texttt{point\_id} chỉ dùng để nhận diện điểm, không được đưa vào phép tính khoảng cách.

\textbf{Luồng giao dịch thẻ 5D.} Tập ban đầu \(D_0\) có 500 giao dịch. Mỗi giao dịch có năm thuộc tính số:

\begin{table}[H]
\centering
\caption{Các thuộc tính của luồng giao dịch thẻ}
\label{tab:card_features}
\normalsize
\begin{tabularx}{\textwidth}{@{}P{5cm}YP{1.7cm}@{}}
\toprule
\textbf{Thuộc tính} & \textbf{Ý nghĩa mô phỏng} & \textbf{Đơn vị} \\
\midrule
\texttt{amount\_vnd} & Giá trị giao dịch & VNĐ \\
\texttt{hour\_of\_day} & Giờ thực hiện giao dịch trong ngày & Giờ \\
\texttt{distance\_from\_home\_km} & Khoảng cách từ nơi giao dịch đến nhà & km \\
\texttt{txns\_last\_7d} & Số giao dịch trong bảy ngày trước đó & Lần \\
\texttt{merchant\_risk\_score} & Điểm rủi ro mô phỏng, từ 0 đến 1 & Điểm \\
\bottomrule
\end{tabularx}
\end{table}

Cột \texttt{point\_id} là mã định danh, không phải đặc trưng để gom cụm. Các giá trị được tạo để mô phỏng, nên không thể xem là mô tả hành vi của khách hàng ngoài đời.

\textbf{Các ví dụ cập nhật và ca kiểm thử bổ sung} 

Tám ví dụ nhỏ minh họa Noise, Creation, Absorption, Merge, Removal, Reduction, Split và Non-split. Mỗi ví dụ có dữ liệu ban đầu, một thao tác cập nhật và tệp \texttt{expected.json} ghi kết quả mong đợi. Một số test dùng tọa độ ngắn hoặc fixture riêng để kiểm tra tình huống biên; đây là các ca chức năng nhỏ, không phải bộ dữ liệu dùng để báo cáo benchmark.

\section{Tập ban đầu và các thao tác cập nhật}

Tập \(D_0\) được gom cụm trước bằng DBSCAN. Sau đó chương trình đọc lần lượt các dòng trong \texttt{updates.csv}: thao tác INSERT thêm một điểm, thao tác DELETE xóa một điểm đang có. File có tổng cộng 300 thao tác, gồm 150 INSERT và 150 DELETE xen kẽ; vì vậy số điểm trong luồng dao động từ 500 đến 501.

Thư mục \path{data/raw/card_stream/} chứa tập ban đầu, luồng cập nhật và file tham số. Trong \texttt{updates.csv}, cột \texttt{ambiguous\_border\_ids} là metadata tùy chọn cho bước đối chiếu các điểm biên; thuật toán không dùng cột này để quyết định cụm.

\section{Các chiều dữ liệu trong benchmark}

Benchmark 2D, 3D và 4D lấy các đặc trưng từ cùng luồng thẻ 5D, không phải ba bộ dữ liệu độc lập. Cấu hình 2D dùng số tiền đã lấy log và giờ giao dịch; 3D thêm khoảng cách từ nhà; 4D thêm số giao dịch trong bảy ngày. Điểm rủi ro của cửa hàng có trong dữ liệu gốc nhưng không dùng trong ba cấu hình này.

Các đặc trưng của thẻ được chuẩn hóa bằng những hằng số cố định lưu trong \texttt{params.json}, rồi dùng cùng dữ liệu đã chuyển đổi cho các thuật toán được so sánh. Hằng số được tính từ toàn bộ luồng mô phỏng để giữ cách biến đổi nhất quán giữa các lượt; đây là giới hạn của thí nghiệm, vì khi chạy online thực tế ta thường chưa biết các giao dịch trong tương lai.

Luồng có 300 cập nhật, nhưng benchmark thời gian giữa bốn thuật toán chạy trên 20 cập nhật đầu để giữ thời gian thực nghiệm vừa phải. Vì 2D, 3D và 4D dùng cùng các điểm và cùng thứ tự cập nhật, kết quả giữa các chiều cho biết việc thêm đặc trưng ảnh hưởng thế nào trong thí nghiệm này; không nên xem chúng như ba lần thử độc lập.

\section{Giới hạn của dữ liệu}

Các bộ dữ liệu chỉ là mô phỏng và có quy mô nhỏ, nên kết quả dùng để minh họa và kiểm tra cách thuật toán hoạt động. Chúng chưa đủ để kết luận thuật toán sẽ có hiệu quả tương tự trên dữ liệu ngân hàng thật hoặc ở quy mô lớn.
